\section{Deriving the column laws and isolating canonical evaluation inputs}
\label{sec:column-laws}

\status{TL1--TL13 reconstruct the actual implication
\texttt{TC.Ok}$\Rightarrow$\texttt{TC.Laws}, including its lattice
finiteness, coefficient, level, and period suppliers. These proofs use the
literal arithmetic and local-index clauses of \texttt{Ok}; they do not
use a previously assembled \texttt{Forms0} or a terminal comparison to
prove its own law hypothesis. TL14--TL21 reconstruct the finite-field,
orientation, and finite-generation parts of the canonical evaluation
constructor, with the remaining local nonvanishing and analytic
algebraicity suppliers identified explicitly. No assertion here produces
the five independent stages of \texttt{Later}, and no Lean build or
kernel verification is asserted.}

\paragraph{Conventions.}
Write $E$ for the fixed number field with involution, $\mathcal O$ for
its integers, and $h(x,y)=\bar x^{\mathsf T}y$ on $E^n$. The field law
used below is
\[
 \tau(\bar a)=\overline{\tau(a)}
 \quad\text{for every embedding }\tau:E\hookrightarrow\mathbb C.
 \tag{TL.1}
\]
For coefficient field $K$, embeddings $e_\tau:E\to K$, and
$\sigma:K\to\mathbb C$, \texttt{FOk} also supplies
$e_\tau(\bar a)=e_\tau(a)^*$ and
$\sigma(b^*)=\overline{\sigma(b)}$. It does not assert the latter
identity for every embedding of an arbitrary larger $K$.
The rank, prime, character, and index data called \texttt{TC.Ok}
throughout this section are the actual source data, including
\texttt{ChIdx.Ok}. Their existence is a separate question, treated
only to the extent stated in TL14--TL21.

\paragraph{TL1: finitely many self-dual Hermitian lattice classes.}
Put $d=[E:\mathbb Q]$, $D=nd$, and $q=|\operatorname{disc}(E)|$.
For a full self-dual $\mathcal O$-lattice $N=N^\#$, its underlying
integer lattice has the positive integral scalar product
\[
 B(x,y)=\operatorname{Tr}_{E/\mathbb Q}h(x,y).
 \tag{TL.2}
\]
Symmetry follows by conjugating the trace, and positivity follows by
writing the trace as the sum of $|\tau(x_i)|^2$. Let $N^*$ denote the
dual for $B$. If $y\in N^*$, then, for all $a\in\mathcal O$ and
$x\in N$, $\operatorname{Tr}(a h(x,y))$ is integral: apply the defining
condition to $\bar a x$. Thus $h(x,y)$ belongs to the codifferent.
The inverse of an integral trace matrix has denominator dividing its
determinant, so $q\mathfrak D^{-1}\subseteq\mathcal O$. Consequently
\[
 N\subseteq N^*,\qquad qN^*\subseteq N,\qquad
 1\le\det B|_N=[N^*:N]\le q^D.                 \tag{TL.3}
\]
In particular nonzero lattice vectors have length at least one and the
covolume is at most $C=q^{D/2}$.

Here is an elementary uniform bound for independent short vectors.
Let $v_D$ be the volume of the Euclidean unit ball and choose
$R\ge1$ with $v_D R>3^{D-1}C$. If the vectors of $N$ of length at
most $2R$ spanned a subspace of dimension $r<D$, then all the lattice
points in any translate of a radius-$R$ ball would lie in one affine
$r$-plane: differences of two such points have length at most $2R$.
Their disjoint radius-$1/2$ balls in that plane give at most
$(2R+1)^r\le3^{D-1}R^{D-1}$ points. On the other hand, averaging the
number of lattice points over a fundamental parallelepiped gives
$v_D R^D/\operatorname{covol}(N)\ge v_D R^D/C$, a contradiction.
There are therefore $D$ independent vectors $v_i$ of length at most
$2R$.

Set $N_0=\sum_i\mathbb Z v_i$. Its integral Gram entries are bounded
by $4R^2$, and $b=[N:N_0]\le(2R)^D$, since
$\operatorname{covol}(N)\ge1$. Moreover
$N_0\subseteq N\subseteq b^{-1}N_0$. For each of the finitely many
Gram matrices and integers $b$, there are finitely many intermediate
groups, as they correspond to subgroups of
$b^{-1}N_0/N_0$.

One must also retain the $\mathcal O$-action, rather than just the
integer quadratic form. Fix an integral basis $a_s$ of $\mathcal O$.
Multiplication by $a_s$ has operator norm at most
$\max_\tau|\tau(a_s)|$ for (TL.2). In the $v_i$ coordinates its
entries have denominator dividing $b$; they are bounded because the
Gram entries are bounded and the inverse Gram matrix is bounded by
the cofactor formula and $\det\operatorname{Gram}(v_i)\ge1$.
There are consequently only finitely many possible multiplication
matrices. Identical Gram, intermediate-group, and multiplication data
give an integer isometry commuting with $\mathcal O$, hence an
$E$-linear isometry. It preserves $h$ because all the identities
$\operatorname{Tr}(a h(x,y))$ agree and the field trace pairing is
nondegenerate. This proves finiteness up to the actual rational
unitary group. For $n=0$ there is just the zero lattice.
Finally the automorphism group of any one lattice is finite: images
of a fixed integer basis lie in finitely many bounded lattice sets.

\paragraph{TL2: decorations, local frames, and coupling cosets preserve
finiteness.}
A valid decoration consists of $N=N^\#$, a nonzero residue-linear
functional at each type prime, and a residue basis at each auxiliary
prime. Values on finitely many $\mathcal O$-generators of $N$
determine these maps on $N$. They determine the localized maps by
clearing a denominator outside the relevant prime; the prescribed
extension is zero outside that localization. All their possible
values lie in finite residue fields. Thus there are finitely many
decorations over each fixed $N$, and TL1 gives finitely many decorated
orbits.

For a fixed decoration, rational frames for $N_P$ form a torsor under
$\operatorname{GL}_n(\mathcal O_{(P)})$. Restricting the permitted
right-frame transformations to a subgroup of finite index leaves
finitely many frame classes. Taking the product over the finitely
many twist primes preserves finiteness. At the coupling step the
additional coordinate is a coset subject to finitely many bounds
$q_P^{a_P}N_P$ modulo $q_P^{b_P}N_P$. Its possible local cosets are
finite: if $a_P<b_P$, this is a quotient of a free local module by
a nonzero power of its maximal ideal, and otherwise there is only
one coset. The map from global coupling cosets to their family of
local cosets is injective by the definition of the simultaneous
local-condition module. Zero supplies a nonempty fiber. At other
steps the coupling coordinate is the fixed empty set.

The elementary orbit argument being used is this: an equivariant
map with finitely many base orbits and finite fibers above chosen
base representatives has finitely many source orbits. Move a source
point over a chosen representative, then choose within that finite
fiber. No finite cardinality of the full rational matrix group is
assumed.

A stabilizer maps injectively into the rational unitary automorphisms
of its decorated lattice. Indeed, once that rational matrix is known,
surjectivity of each auxiliary residue basis determines the auxiliary
matrix, and equality of each rational local frame determines its
right-frame matrix. Coupling coordinates add a condition, not a new
kernel. Stabilizers are therefore finite. Their orders divide the
corresponding lattice-automorphism orders. The finite set of those
orders, over the finite rank set in \texttt{gamRanks}, has a positive
least common multiple \texttt{DGam}.

\paragraph{TL3: the actual twist levels have finite index and finite
characters.}
For any nonzero integer $c$, the quotient
$\mathcal O_{(P)}/c\mathcal O_{(P)}$ is finite: it is a localization
of the finite ring $\mathcal O/c\mathcal O$. A subgroup of
$\operatorname{GL}_n(\mathcal O_{(P)})$ containing the principal
congruence kernel modulo $c$ therefore has finite index.

For the actual intermediate index matrix $h$ at step $j$, the
\texttt{boundQ} clause of \texttt{ChIdx.Ok} says
\[
 h^{-1}\Lambda\subseteq\Lambda,
 \qquad q^j h\Lambda\subseteq\Lambda.             \tag{TL.4}
\]
If $\kappa\equiv1\pmod{q^{j+1}}$ on the first coordinate block,
put $a=h^{-1}\operatorname{diag}(\kappa,1)h$. For $x\in\Lambda$,
the first coordinates of $q^jh x$ are integral, so
$(\kappa-1)\operatorname{ini}(h x)\in q\mathcal O_{(P)}^n$.
The first inclusion in (TL.4), applied to the embedded first block,
gives $(a-1)\Lambda\subseteq q\Lambda$. The same calculation for
$\kappa^{-1}$ proves $a^{-1}\Lambda\subseteq\Lambda$ and hence
$a\Lambda=\Lambda$. Its reduction is the identity on every
specified residue flag. This is precisely membership in the source's
\texttt{Jat}, whose first-block projection defines \texttt{levQ}.
Thus the actual level contains the congruence kernel of depth $j+1$.
At the top step the level is generated by an extended row level and
the congruence kernel of its specified depth plus one, so the same
conclusion holds. The row level itself is a preimage of a subgroup
of the finite residue general linear group.

The level characters are products of finite residue multiplicative
characters, finite residue additive characters, and the prescribed
local additive character on $q^{-D}\mathcal O_{(P)}$ modulo
$\mathcal O_{(P)}$. Every value is a root of unity: these additive
quotients and residue multiplicative groups are finite. Products and
inverses retain that property. The local-index character formula in
\texttt{ChIdx.Ok} identifies these products with the actual level
characters. This argument proves their finite order after the index
data have been supplied; it does not construct those index data.

Combining TL2 and these facts gives all seven clauses of
\texttt{GStep.Ok}: finite orbits, finite stabilizers, unitary rational
matrices, root-of-unity characters, exact transport of $L^\#$, fullness
of that lattice, and stabilizer injectivity. Replacing the preferred
orbit representative changes none of these seven assertions.

\paragraph{TL4: the zero step is free and all its twists have the same
classes.}
At step zero the auxiliary level has unipotent reduction at a prime
above $g_0$, and the base prime bound is
\[
 [E:\mathbb Q]\,k=2fk<g_0-1.                     \tag{TL.5}
\]
Let $\gamma$ stabilize a decorated point. By TL2 it has finite order.
Its prime-to-$g_0$ part is congruent to the identity: a unipotent
matrix in residue characteristic $g_0$ has $g_0$-power order. A
finite-order matrix $u=1+X$ with order $r$ prime to $g_0$ and
$X\in P$ is the identity. Indeed $(1+X)^r=1$ first puts $X$ in
$P^2$, and iteration puts it in every $P^a$, whose intersection is
zero. Hence $\gamma$ has $g_0$-power order. If it were nontrivial,
a power would have order $g_0$ and its rational characteristic
polynomial on the underlying $[E:\mathbb Q]k$-dimensional space
would contain $\Phi_{g_0}$, of degree $g_0-1$. This contradicts
(TL.5). Stabilizer injectivity now kills the auxiliary and local
frame matrices as well. The zero-step action is free.

The actual index at zero is unique by \texttt{zero\_uniq}; its
auxiliary level is the same \texttt{colLev 0} for every twist, its
twist-frame level is the full local general linear group, and its
coupling coordinate is empty. Forgetting those local frames
therefore gives a bijection with the same decorated orbit set for
every twist. To see injectivity directly, any two frames over an
identical local lattice differ by an allowed integral invertible
matrix. Adding frames proves surjectivity. The prescribed preferred
representative minimizes the same decorated-orbit datum, so the
preferred decorations agree literally. This proves the actual
\texttt{cl0} bijections and their decoration equality.
Finally \texttt{tup mu0}=1 and the injective finite-idele reading
identity imply $\xi_{\mu_0}=1$. Every other read character has finite
order, hence has complex absolute value one. The defining products
for $\rho$ and $\rho'$ are therefore of modulus one. TL2's positive
\texttt{DGam}, its divisibility, and the trivial zero stabilizers
give the complete chain laws.

\paragraph{TL5: covariance and isometry of the coefficient family.}
This paragraph starts from a generic \texttt{GStep.Ok}, independently
of the final column-law record. Let $\lambda$ be any supplied
dominant tensor weight, $V_\lambda$ its polynomial coefficient space,
and
\[
 R(g)=\chi(g)\,R_\lambda(e_\tau(\gamma(g))) .       \tag{TL.6}
\]
Exact lattice transport says
$L^\#(gd)=\gamma(g)L^\#(d)$. Extension along $e_\tau$ carries this
to equality of coefficient lattices; taking coordinate duals changes
the frame by inverse transpose. The defining minor polynomials and
their divided-power translates then give exactly
\[
 \nabla(gd)=R(g)\nabla(d),\qquad
 \Delta(gd)=R(g)\Delta(d).                         \tag{TL.7}
\]
Here the character scalar causes no integral error. If
$\chi(g)^r=1$, both $\chi(g)$ and its inverse are algebraic integers,
so multiplication by it preserves every $\mathcal O_K$-module.

The rational matrix is unitary. Under every $e_\tau$ this is the
unitary identity for the coefficient pairing, since \texttt{FOk}
commutes with conjugation on $E$. The normalized tensor Fock form
is invariant under unitary changes of variables; this follows on
monomials from the tensor symmetric-power scalar product and then
on the Weyl subspace by restriction. At the distinguished embedding,
a root of unity has absolute value one. Injectivity of $\sigma$
and \texttt{FOk.sig\_star} give
$\chi(g)^*\chi(g)=1$. Thus (TL.6) also preserves the pairing.
Together with the finite orbit and stabilizer clauses these are all
the assertions of the actual \texttt{fam.Ok}.

\paragraph{TL6: rational fixed vectors give a basis of the complex
forms.}
Choose one representative of each class $c$ and write $\Gamma_c$
for its finite stabilizer. The rational forms are precisely
$\prod_c V_\lambda^{\Gamma_c}$. The matrix coefficient realization
$\operatorname{ev}:V_\lambda\to\operatorname{Val}_{\mathbb C}$ is
injective and intertwines $R(g)$ with
$F(Y)\mapsto\sigma(\chi(g))F(\gamma(g)^{-1}Y)$; evaluation of the
minor generators proves this identity. In a fixed polynomial basis
its columns remain independent after any field embedding: an
injective finite matrix has a left inverse, which can be embedded.
Thus rationally independent form vectors have complexly independent
evaluations. The defined complex polynomial value space is the
complex span of the evaluated standard coefficient basis (equivalently
the complex matrix-coefficient realization of the same Weyl module).
Thus evaluation after scalar extension is onto this space as well
as injective; no surjectivity onto a larger function space is assumed.

For spanning, the projector
\[
 e_c=|\Gamma_c|^{-1}\sum_{g\in\Gamma_c}R(g)        \tag{TL.8}
\]
is a matrix over $K$. Its image is exactly the invariant subspace:
translation permutes the summands, and it fixes an invariant vector.
The same identities hold after scalar extension. Hence every complex
stabilizer-fixed value is a complex linear combination of rational
fixed values. Apply this separately at the finitely many classes.
The actual chosen rational basis therefore evaluates to a basis of
\texttt{Fi}, proving \texttt{LawBasis} and finite dimensionality.
Division in (TL.8) takes place only in the characteristic-zero
field; it is not an assertion of integral averaging.

\paragraph{TL7: the full integral forms lattice and the actual
all-place Hermitian metric.}
For a full lattice $M\subset E^n$, its extension along each embedding
and its coordinate dual are full. The corresponding integral Weyl
lattices are full and finitely generated. One direct verification
puts $M$ between two nonzero scalar multiples of the standard
lattice. The finite degrees of the minors and determinant twists
put its Weyl lattice between two scalar multiples of the standard
Weyl lattice, whose finite divided-power construction spans the
coefficient representation. The same argument applies to the
$\nabla$ lattice and its dual $\Delta$.

For finitely many classes and full coefficient lattices $M_c$, put
\[
 \mathcal L=\{x\in\textstyle\prod_cV_\lambda^{\Gamma_c}:
                         x_c\in M_c\ \forall c\}. \tag{TL.9}
\]
Its inclusion in $\prod_c M_c$ makes it finitely generated, since
$\mathcal O_K$ is Noetherian. For any rational form $x$, choose
nonzero $a_c\in\mathcal O_K$ with $a_cx_c\in M_c$. The product
$a=\prod_c a_c$ puts $ax$ in $\mathcal L$ and is nonzero.
Thus $K\mathcal L$ is the whole forms space.

The all-place positivity needs care about the chosen basis. Define
\[
 K^{\rm cm}=\{b\in K:\phi(b^*)=\overline{\phi(b)}
                       \text{ for every }\phi:K\to\mathbb C\}.
 \tag{TL.10}
\]
This is a subfield. It contains every $e_\tau(E)$ by (TL.1), and
every character root of unity: first use \texttt{sig\_star} and
injectivity at $\sigma$ to get $b^*=b^{-1}$, then use $|\phi(b)|=1$
at any $\phi$. All representation matrices (TL.6) therefore have
entries in $K^{\rm cm}$. Row reduction of the finite invariant
equations produces a basis with coefficients in that subfield.
The source's \texttt{CSD.basis} selects precisely such a basis when
this rational-basis condition holds. This is not an assertion that
an arbitrary $K$-basis has conjugation-compatible coefficients.

Let $A_c$ be the matrix of these basis values at $c$, and $G$ the
standard normalized tensor Fock Gram matrix. The actual class-sum
matrix is $H=\sum_c |\Gamma_c|^{-1}A_c^*GA_c$. Because only the
basis coefficients and the fixed representation pairing enter,
(TL.10) gives for every embedding $\phi$
\[
 \phi(H)=\sum_c |\Gamma_c|^{-1}
               \phi(A_c)^*\phi(G)\phi(A_c).       \tag{TL.11}
\]
The matrix $\phi(G)$ is positive definite: the standard Fock
monomial scalar product is positive and restriction to the Weyl
space is positive. The stacked matrices $A_c$ have trivial kernel;
a left inverse shows that they still do so after $\phi$.
Consequently every nonzero vector has a strictly positive summand
in (TL.11). This proves positivity at every infinite place and,
together with (TL.9), the literal \texttt{LawGood}.

At $\sigma$ the evaluation isometry identifies that class sum with
\[
 \|x\|^2=\sum_c\frac{1}{|\Gamma_c|}
               \operatorname{ipVal}(\operatorname{ev}x_c,
                                     \operatorname{ev}x_c).
 \tag{TL.12}
\]
The normalizing denominator is the product of the Fock
self-pairings of the highest vectors, hence of their squared
norms. It is not the product of their norms. The source's infinite
place representing $\sigma$ may choose its conjugate embedding;
the real value of the self-pairing is unchanged. The positive square
root frame of $\phi(H)$ consequently gives exactly the source's
\texttt{majAt} and \texttt{LawNorm}, with its original normalization.

\paragraph{TL8: the residue ranks, distinguished vector, and dilation
laws.}
For a full rank-$n$ lattice $M$, localization at $P$ is free of rank
$n$ over the discrete valuation ring $\mathcal O_{(P)}$. The natural
map $M/PM\to M_P/PM_P$ is an isomorphism: a denominator outside
$P$ is invertible modulo $P$, and if $sx\in PM$ with $s\notin P$,
choose $t$ with $st\equiv1\pmod P$ to recover $x\in PM$.
Both residue modules in \texttt{TDat} therefore have dimension $n$.
The supplied additive, residue-linear $f$ kills $PN$, so it factors
uniquely to a nonzero linear functional $\mathrm{fv}:N/PN\to k_P$.
Choose $e\in N$ with $f(e)=1$, by scaling a nonzero value by a lift
of its residue inverse. The clause \texttt{te\_mem} gives
$Te\in\mathrm{dn}N\setminus P\mathrm{dn}N$, so its residue class
$\mathrm{te}$ is nonzero. TL11 verifies that clause for the actual
decorated lattices.

Dilation by $a\in k_P^\times$ replaces only $f_P$ by $af_P$.
Its inverse is dilation by $a^{-1}$, it commutes with the acting
group, and it preserves $N$ and $L^\#$. It therefore induces a
group action on classes. The raw stabilizer is unchanged; changing
to the preferred representative conjugates it, so its cardinality
is also unchanged. If $g$ takes the preferred representative of the
dilated class to the dilation of the old representative, its
rational matrix induces an isomorphism
$e:N_{c_a}/PN_{c_a}\to N_c/PN_c$ and
\[
 \mathrm{fv}_{c_a}(v)=a\,\mathrm{fv}_c(e(v)).      \tag{TL.13}
\]
This proves all the literal \texttt{TDat.Laws}, with the factor
$a$ in its correct direction.

\paragraph{TL9: the diamond action preserves complex and integral
zero-step forms.}
Since the zero-step action is free, there is a unique group element
$g_d$ carrying the preferred representative of $[d]$ to $d$.
It satisfies $g_{hd}=hg_d$ and $g_{\operatorname{rep}c}=1$.
Extend a class value $F_c$ to points by
$\widetilde F(d)=\operatorname{tr}(g_d)F_{[d]}$.
Commutation of dilation with the group action proves directly that
pullback along dilation has identity and multiplication laws. The
residue-unit group is commutative, so the order reversal for
pullbacks causes no change to the source's formula. The evaluation
intertwining in TL6 shows that it preserves the polynomial value
space; stabilizer invariance is automatic at this free step.

For integral values, put
$y_c=R(g_{\operatorname{dia}_a\operatorname{rep}c})
 x_{[\operatorname{dia}_a\operatorname{rep}c]}$.
Equation (TL.7), and the fact that dilation leaves $L^\#$ unchanged,
put $y_c$ in the actual coefficient lattice at $c$. It is a rational
form by freeness. Its evaluation is the diamond transform of the
evaluation of $x$. This proves the precise existential
\texttt{dia\_lat} statement, in addition to complex stability and
the representation laws.

\paragraph{TL10: spanning at every coefficient place, and
\texttt{laws\_all}.}
The same proof as (TL.9) applies to the coefficient lattice built
from \texttt{ker} at any step and selected prime: \texttt{TDat.Ok}
supplies its fullness. Thus the actual \texttt{BLat} spans the
rational forms space. For any coefficient place $w$, its localization
$\operatorname{loc}_w(\texttt{BLat})$ contains \texttt{BLat}, so it
has the same full $K$-span. This holds for every place in the supplied
\texttt{Places} record; no seen-place or residue-characteristic
restriction is introduced.

Now start with arbitrary number fields and a generic chain
$I:\texttt{ChainIn}$, level datum $J:\texttt{AtIn}$, period datum
$Q:\texttt{PerIn}$, coefficient embeddings, dominant weights, and
places. Under exactly \texttt{FOk}, \texttt{I.Ok}, \texttt{J.Ok},
and \texttt{Q.Ok}, TL5--TL7 apply at every chain step and also to
$Q.S'$ with its shifted terminal weight. TL8 supplies every level
law, TL9 supplies the zero-step diamond laws and integral stability,
and the preceding denominator argument supplies every
\texttt{B\_span}. These are precisely all fields of
\texttt{J.Laws} and \texttt{FmLaws}. In particular the genuine
proof of \texttt{laws\_all} requires no terminal incidence or
polynomial-value comparison premise.

\paragraph{TL11: the actual level datum is valid and tied to the
decorated lattices.}
We now return to \texttt{TC.Ok}. Its base prime laws give a rational
prime $T$, selected primes $S_T$ with $|S_T|=f$ and $|k_P|=T$,
no conjugate pair in $S_T$, and $T\in P\setminus P^2$.
Auxiliary and twist-frame primes avoid $P,\bar P$: the former have
different rational residue characteristics, and the tuple excludes
ramification at $T$. Thus the source's \texttt{MoveHyp} holds.

Here is the local lattice calculation needed for the actual data.
The elementary duality formula NB1 depends only on lattices and
LP1, and is used in that form. If $N=N^\#$ and
$K_P^{\rm gl}=\{x\in N:f_P(x)=0\}$, then the down lattice
\[
 D=K_P^{\rm gl}+P(K_P^{\rm gl})^\#                \tag{TL.14}
\]
has localizations $D_P=\ker(f_P|_{N_P})$,
$D_{\bar P}=((K_P^{\rm gl})^\#)_{\bar P}$, and $D_R=N_R$ elsewhere.
The exceptional two are dual to each other, proving $D=D^\#$.
Choose a local basis with $f_P$ its last residue coordinate. Then
$D_P=\mathcal O_{(P)}^{n-1}\oplus T\mathcal O_{(P)}$ in that basis.
If $f_P(e)=1$, the last coordinate of $Te$ is a unit times $T$;
it lies in $D_P$ but not in $PD_P$. At $\bar P$ the down lattice
contains $N_{\bar P}$, and at all other primes it equals $N_R$.
Hence $Te\in D\setminus PD$ globally. This proves the actual
\texttt{te\_mem}.

For the up operation with $f_P(x)=1$, put
$A=N+P^{-1}x$ and $U=A^\#+\bar P A$. Its localizations are
$U_P=N_P+T^{-1}\mathcal O_{(P)}x$,
$U_{\bar P}=(A^\#)_{\bar P}$, and $U_R=N_R$ elsewhere.
Again the two exceptional lattices are dual and $U=U^\#$.
The updated functional is $f_P^+(y)=f_P(Ty)$ on $U_P$; in the
displayed basis its kernel is $N_P$. Given the down move's specified
nonzero residue-linear map $g:D/PD\to k_P$, its actual functional is
$f_P^-(y)=\rho_P(s)^{-1}g([sy])$ for $y\in D_P$, where
$s\notin P$ and $sy\in D$, and is zero outside $D_P$.
More generally, any ambient down functional giving a valid decoration
descends to such a map on $D/PD$, and its localization and off-lattice
clauses force this same formula.
Replacing $s$ by another denominator does not change the value:
multiply the two numerator equations and use residue-linearity.
For down after the specified up, (TL.14) returns $D=N$ and the
specified $g$ is the original residue map induced by $f_P$.
The displayed formula is then $\rho_P(s)^{-1}f_P(sy)=f_P(y)$.
Thus both the lattice and its functional return literally to the
old ones. If $x-x'\in PN$, then
$P^{-1}(x-x')\subset P^{-1}(PN)\subset N$, so
$N+P^{-1}x=N+P^{-1}x'$ globally. Their raw up lattices and
functionals are therefore equal even for invalid inputs. Thus both
the validity and fallback branches agree, and up depends only on
the stated residue class without an additional validity assumption.

All remaining decorations and local frames are unchanged, since
their primes avoid the exceptional pair. The coupling coordinate is
defined by local conditions at those other primes, so its allowed
set is unchanged as well. Thus these are literal operations on the
source points, not only operations on lattice isomorphism classes.
Applying a rational unitary matrix transports sums, duals, and
residue maps; hence the up/down action identities hold exactly,
including their frame and coupling coordinates. Dilation simply
rescales one functional, proving all its raw identities. Fullness
of the level kernel follows, for example, by placing it between a
nonzero integer multiple of $L^\#$ and $L^\#$; the finite residue
conditions provide such an integer.

These calculations establish all \texttt{TDat.Ok} and
\texttt{TDat.Tied} clauses for \texttt{tdatAt}. Rechoosing preferred
representatives changes none of them. The Weil additive character
is nontrivial because the specified $\xi_0$ is a unit at the selected
prime; its value on a generator of the inverse-prime quotient has
nontrivial residue phase. Together with the base prime/cardinality
laws, TL4's zero-step freeness, and the positive constants already
in the base data, this gives the literal \texttt{atIn.Ok} for every
word and every twist. This proof does not use NB7/NB8 or the final
\texttt{Forms0} constructor.

\paragraph{TL12: the actual period datum and its level-prime overlap.}
Let $n=N_m$, $n'=n-1$, and $\pi=\pi_T$ with
$\pi\bar\pi=T$. The row step has no type primes. TL2--TL3 apply to
its residue row level and its reduced twist level, giving
\texttt{stepX.Ok}. Extend its rational, auxiliary, and twist acting
matrices by a last coordinate equal to one. This gives an
injective acting-group homomorphism: its first block recovers the
original element. Its rational matrix is
$\operatorname{diag}(\gamma,1)$, so the first coordinates transform
by $\gamma$ and the last coordinate stays fixed.
The point's twist frame, however, extends as
$\operatorname{diag}(\mathrm{fr},\pi)$. Its exact equivariance is
\[
 \operatorname{diag}(\gamma,1)
 \operatorname{diag}(\mathrm{fr},\pi)
 \operatorname{diag}(q,1)^{-1}
 =\operatorname{diag}(\gamma\mathrm{fr}q^{-1},\pi).
\]

The top auxiliary level is exactly the image of the row level under
this diagonal embedding. At a twist prime the top level is its
extended row level times the specified principal congruence group.
The character is pulled back from reduction, is trivial on that
congruence group, and on $\operatorname{diag}(g,1)$ equals the row
character: every old block determinant is unchanged and the added
block determinant is one. These facts prove the exact
\texttt{chr\_ext} identity and the equivariance of the extended
framed decoration. Since $T$ is excluded at every twist prime,
$\pi$ is a unit there, so the specified extended frames are valid.

For a row lattice $N=N^\#$ the extended decoration has
\[
 \widetilde N=N\oplus\mathcal O,\qquad
 \widetilde f_P(x,z)=\rho_P(\bar\pi z),\qquad
 \widetilde L=N\oplus\pi\mathcal O,
 \qquad\widetilde L^\#=N\oplus\bar\pi^{-1}\mathcal O.
 \tag{TL.15}
\]
To check the third equality, the selected primes are exactly the
prime divisors of $(\pi)$ with exponent one: its norm is $T^f$ and
the $f$ selected primes each have norm $T$. Their intersection is
therefore $(\pi)$, and $\bar\pi$ is a unit at each one. The last
equality follows by taking the Hermitian dual; in particular the
denominator is $\bar\pi$, not $\pi$.
With $t_z=\bar\pi$, (TL.15) says exactly
\[
 (x,z)\in\widetilde L^\#\iff x\in L^\#\text{ and }t_zz\in\mathcal O,
 \quad (x,z)\in\widetilde N\iff x\in L^\#\text{ and }z\in\mathcal O,
 \quad \bar t_z t_z=T.                            \tag{TL.16}
\]

The terminal preferred-representative predicate asks to lie in the
image of this extension. Every extended point itself satisfies it,
so the representative selected in its orbit also lies in that image.
This proves the actual \texttt{rep\_til}; it is not a claim that
all terminal orbits lie in the image. The subgroup $H$ in this
period datum is the full row acting group, so its finite index is
one and its rational-image condition is immediate.
The preceding identities prove every field of \texttt{perIn.Ok}
and \texttt{perIn.Over(atIn)}, for every word and twist, directly
from \texttt{TC.Ok}. No terminal comparison from TG is used.

\paragraph{TL13: the complete implication \texttt{Ok}$\Rightarrow$
\texttt{Laws}, and its exact scope.}
For any actual base $B$, prechoices $X$, and true-column data
$I:\texttt{TC}(B,X)$, assume only $I.\texttt{Ok}$ in addition to the
data already built into $B,X,I$. TL3--TL4 give all \texttt{ch} fields
and nonempty zero indices; TL11 gives \texttt{lev}; TL12 gives
\texttt{per} and \texttt{ovr}; TL10 gives \texttt{atL} and
\texttt{fmL} at every word and twist. The base fields and primes
give \texttt{pr}, TL4 gives \texttt{xi0} and \texttt{bij0}, and
every decorated level kernel is a full lattice by the finite-residue
argument in TL11, giving \texttt{lat}. These are all eleven fields
of the literal \texttt{TC.Laws} record.

Thus the law premise in the previously conditional constructions
can now be supplied by this proof whenever \texttt{TC.Ok} is
supplied. The quantifiers are all words, all character indices,
all natural-number chain steps, all classes, and, in
\texttt{B\_span}, all coefficient places. No new \texttt{Pins},
\texttt{Later}, chosen adelic realizer, or terminal comparison is
needed to derive these laws. This implication alone does not supply
\texttt{TC.Ok}; in particular its local index nonvanishing clauses
remain present in that input.

\paragraph{TL14: one number field contains all the normalized frame
character values.}
We give the substantive arithmetic argument behind
\texttt{frame\_values}. Let $\xi$ be a continuous Hecke character
of a CM field, with infinity type $n_\tau\in\mathbb Z$, and let
$r\in\mathbb Z$ satisfy $n_\tau-r\in2\mathbb Z$ for every CM
embedding. Put
\[
 a_P=\xi(P)\,\mathrm N(P)^{-r/2}
 \quad\text{at unramified primes }P.              \tag{TL.17}
\]
The complex target has a neighborhood of one containing no
nontrivial subgroup. Continuity on the finite ideles therefore
gives a finite set $S$ outside which local units are killed and,
at each $P\in S$, an open principal-unit subgroup in the kernel.
One can see this without a global character theorem: choose a
basic open subgroup inside the inverse image of that neighborhood.
Its image is a subgroup contained in the neighborhood, hence is
trivial. The quotient of local units by the resulting congruence
subgroup is finite. Thus the local character values on units at
$S$ have bounded root-of-unity orders.

Let $M$ contain all conjugates of $E$ and those finitely many roots
of unity. The ordinary ideal class group is finite. For clarity,
its usual lattice proof only needs a uniform short-vector bound:
for an integral ideal $I$, averaging a ball of volume exceeding
its Minkowski covolume gives a translated ball containing two
lattice points. Their difference yields $0\ne a\in I$ with
$|\mathrm N(a)|\le C_E\mathrm N(I)$. Hence
$(a)I^{-1}$ is integral of bounded norm. There are finitely many
ideals of bounded norm, since an ideal of norm $b$ contains
$b\mathcal O$ and corresponds to a subgroup of the finite group
$\mathcal O/b\mathcal O$. This proves class-group finiteness.

Let its order be $h$. If unramified primes $P,P'$ have the same
class, then $P(P')^{h-1}=(a)$ for a nonzero algebraic integer $a$ which is
a unit at every prime of $S$. The Hecke product formula gives
\[
 a_Pa_{P'}^{h-1}
  =\left(\prod_{Q\in S}\xi_Q(a)\right)^{-1}
    \prod_\tau \tau(a)^{-n_\tau}
       \bigl(\tau(a)\overline{\tau(a)}\bigr)^{(n_\tau-r)/2}
  \ \in M.                                         \tag{TL.18}
\]
Indeed $\xi_\infty(a)=\prod_\tau(\tau(a)/|\tau(a)|)^{n_\tau}$,
and $\mathrm N((a))=\prod_\tau|\tau(a)|^2$; the parity assumption
is exactly what turns the remaining powers of absolute values
into the integral powers in (TL.18).
Choose one unramified prime in each class which occurs. Setting
$P=P'$ in (TL.18) makes its chosen value algebraic over $M$.
Adjoin these finitely many nonzero values. Equation (TL.18) now
places every $a_P$ in this one number field. This proves the
literal \texttt{ValuesIn(r/2)} assertion, including its uniformity
over all unramified primes, rather than just individual
algebraicity.

For the supplied frame, $\mu_j$ has even infinity exponents,
$\eta$ has exponent one and uses $r=1$, and $\chi_\Psi$ has
exponent two and uses $r=0$. The finite-order $\delta,\delta'$
and the finite auxiliary characters require only cyclotomic
extensions. A finite compositum, then its normal closure, gives
one finite Galois field containing all five character families,
all field embeddings, and the auxiliary roots, exactly as in
\texttt{frame\_values}. A primitive generator may be required
inside the fixed coefficient field in \texttt{large\_2}; this
does not amount to enlarging that fixed field automatically.

\paragraph{TL15: finite root data, actual residue characters, and
finite-idele readings.}
For a finite-order character $\xi$ of order dividing $N>0$, its
values are powers of a primitive $N$th root. At each of its
finitely many ramified primes with norm $q$, adjoin
$\sqrt q$ and a primitive $q^{2m}$th root. Taking the union for
a finite family $\Xi$ gives finitely many algebraic generators
which imply the exact \texttt{Holds} condition for every member.
This includes the whole character image, not just its values at
uniformizers. A field containing one primitive $N$th root
contains every $N$th root by its cyclic list of powers.

At an auxiliary prime above $g_i$, identify its residue field
with $\mathbb F_{g_i}$ by the prescribed prime-field isomorphism.
Lift the primitive multiplicative character $\omega_i$ and the
additive character $b\mapsto\exp(2\pi\mathrm i b/g_i)$ through
the injective coefficient embedding. Define
\[
 \psi_{0,P}(b)=\psi_i(\overline{\xi_0}\,b),\qquad
 \chi_{j,P}(b)=\omega_i(b)^{2(j+1)}.               \tag{TL.19}
\]
Here $\overline{\xi_0}$ means reduction at $P$, not complex
conjugation. Since $\xi_0$ is a local unit, $\psi_{0,P}$ is
nontrivial. The primitive $\omega_i$ has order $g_i-1$; the
inequality $2m<g_i-1$ makes the powers in (TL.19) distinct.
The powered characters themselves need not be primitive.
The reductions of $\xi_0$ at the selected primes over a given
$g_i$ agree under these prime-field maps: $\xi_0$ lies in the
specified quadratic subfield, and all these primes lie over
the one selected prime of that subfield. Thus the additive and
multiplicative compatibility clauses hold with their actual
residue maps. Tameness makes a local unit character depend only
on its residue; compare a unit to the natural integer with the
same residue and use the prescribed formula for $\mu_j$ on that
integer. This proves the inverse-character tie in
\texttt{Chars.Tied}. Lift the normalized $\mu_j$ and $\eta$
prime values using TL14. The exponential in (TL.19) gives
\texttt{WeilTied} literally. These arguments construct all of
\texttt{Chars.Ok}, \texttt{Tied}, and \texttt{WeilTied} from the
stated frame and root data.

For the reading step, a finite-order character is trivial on
the infinite ideles: every complex unit has an $N$th root, so
its value under an order-$N$ character is one. Its finite-idele
restriction therefore determines it. If all its values lie in
$\sigma(K)$, pull them back uniquely to $K^\times$. Multiplication
and every triviality identity descend through the injective
$\sigma$. Continuity provides the open congruence kernel as in
TL14. The supplied \texttt{IsTwist} conditions say precisely
finite order, triviality on the base-field ideles, tameness, and
ramification only at the selected split, degree-one unramified
primes outside the excluded set. They give the tuple's wild-unit,
inert-place, and base-field kernel clauses. For an infinite
product of local units the assertion follows first on the
finite-support subgroup and then by its density and continuity.
Conversely the tuple's open kernel gives a continuous finite-idele
character, and the rational-unit identity descends it to the
idele class group with trivial infinite component. These two
constructions are inverse. In particular the read tuple is
unique and the tuple for the trivial character is one.

\paragraph{TL16: orienting the auxiliary factors and obtaining the
quotient part of \texttt{Pins}.}
Suppose an integral $a$ lies in the specified imaginary quadratic
subfield and satisfies $a\bar a=g$, with $g$ a rational prime.
Neither $a$ nor $\bar a$ is a unit: conjugation would otherwise
make both units. Their principal ideals in the quadratic field
have norms greater than one with product $g^2$, so each has norm
$g$ and is a maximal ideal. A prescribed prime above $g$ contains
one of them. Choose $b=a$ or $b=\bar a$ so that the ideal of
$\bar b$ is that prescribed prime. A prime of $E$ contains
$\bar b$ if and only if its restriction is that prime. Norm and
the required congruence to one persist under conjugation. Doing
this independently at $g_0,g_1$ proves the actual
\texttt{exists\_oriented} statement.

Under the given $\texttt{Ctx.Over}$ relation replace the type-prime
factors by the context's own $\pi_T(u)$, preserving the now-oriented
auxiliary factors. Their prescribed norm and congruence laws make
this another valid \texttt{PreX}. The context defines
$c(u)=\pi_T(u_0)/\pi_T(u)$, and its prime norm makes the denominator
nonzero. Hence
\[
 c(u)\pi_T(u)=\pi_T(u_0)\quad\text{for every word }u. \tag{TL.20}
\]
This supplies \texttt{Pins.cof}. The other field,
\texttt{Pins.iX}, is the literal agreement of the supplied row
frame with the first columns of the full frame; it is still an
explicit antecedent in \texttt{evLg\_canonical}. No change of
coordinates here proves an arbitrary supplied frame inclusion.

\paragraph{TL17: finite principalization and the precise period-field
boundary.}
For any nonzero ideal $I$ of a number field, class-group finiteness
gives $I^h=(\beta)$ for some $h>0$. Adjoin a root $x^h=\beta$.
It is integral, and in the larger integer ring
$(x)^h=(I\mathcal O')^h$. Unique ideal factorization then gives
$(x)=I\mathcal O'$ by comparing exponents at every prime.
Fractional ideals are treated after clearing a denominator.
Thus finitely many rank-one fractional lattices become principal
over one finite extension, with algebraic generators. Applied to
the coordinate-axis sections of the module generated by the
conjugate vectors, this supplies the algebraic derivative scalars
$c_D$. Their sections commute with this extension: over a
Dedekind domain the larger integer ring is torsion-free, hence
flat, so taking the kernel of the projection to the other
coordinates commutes with extension. This proves the finite
principalization content of \texttt{exists\_dgen}.

For the period field, the host identity
$\mathfrak D_{E/K_0}=\mathfrak b\bar{\mathfrak b}$ gives
$\mathfrak b^{-1}\mathfrak D_{E/K_0}=\bar{\mathfrak b}$.
Hence each source period-line lattice is the extension of a
conjugate of $\bar{\mathfrak b}$. The preceding root construction
principalizes every one simultaneously. Adjoining also
$j(\mathbb C/\mathfrak c)$ therefore gives the actual
\texttt{FieldOK}, provided this CM $j$-value is algebraic.
The source proves the stronger integrality assertion using the
modular polynomial/CM-lattice theorem. That analytic modular
input is retained here; it has not been inferred from ideal
principalization. The finite extension and its generator are
constructed once that one algebraicity input is supplied.

Similarly, a full \texttt{AmbIn} includes the chosen period pair,
the quadratic trace isometry and its complex linear extension,
and the local correction data, as well as $c_D$ and a discriminant
twelfth root. TL17 closes its finite principalization part, not
the separate supplier of these period and trace-isometry data.
These are recorded as exact inputs in TL21.

\paragraph{TL18: reducing normalized-jet algebraicity to the primitive
theta inputs.}
There is a useful elementary step in the analytic supplier which
can be reconstructed without presuming its conclusion. Consider
an algebra of smooth germs over $\overline{\mathbb Q}$, generated
by germs whose derivatives in algebraic coordinate directions
stay in the same algebra and whose values at zero are algebraic.
The sum, product, and chain rules prove by induction that every
iterated derivative of every germ in the algebra is algebraic
at zero. The same holds after a linear substitution with algebraic
coefficients.

Suppose the supplied four-copy identity places
$f(z_1)f(z_2)f(z_3)f(z_4)$ in this algebra. Apply the same
multi-derivative separately in each copy and set all four
variables to zero. The result is $(D^\alpha f(0))^4$, which is
algebraic. Hence $D^\alpha f(0)$ is algebraic too. This argument
works even when $f(0)=0$; it never divides by a value of $f$.
If the available algebraic directions are a spanning family,
choose independent columns and invert their nonsingular minor
over $\overline{\mathbb Q}$ to express any algebraic direction
with algebraic coefficients. This proves the source's direction
descent in the four-copy calculation.

For the actual vacuum, the four-copy identity is a finite sum
of products of one-dimensional torsion translates of the
normalized even theta function. Its coefficients are the
algebraic correction factor $\kappa$, signs, and inverse square
roots of finite indices. Each factor of $\kappa$ satisfies its
stated $2h_{K_0}$th-power equation with an algebraic right-hand
side, so it is algebraic; only finitely many nontrivial factors
occur. Thus, \emph{given} that four-copy identity and the
normalized CM torsion-translate Taylor theorem, the preceding
argument proves algebraicity of the standard vacuum's Taylor
coefficients.

The remaining transport is also explicit. Two rational
Heisenberg translations compose with a phase
$\exp(2\pi\mathrm i q)$, $q\in\mathbb Q$, so their extra factor
is algebraic. The intertwining operator is a finite sum of such
translations with coefficients equal to signs times an inverse
square root of an index. It preserves algebraicity of all
torsion-translate Taylor coefficients. The source's normalized
indicator theta function is, by its explicit finite presentation,
a finite algebraic linear combination of rational translations
of this moved vacuum. The multi-derivative and coordinate
reindexing formulas then prove \texttt{JetAlg}, and the supplied
entirety identities give \texttt{ESmooth}.

The remaining primitive statements are therefore precise: the
normalized CM theta Taylor theorem, its torsion-translation
algebraicity theorem, and the vacuum/four-copy and normalized
indicator identities. This paragraph proves the closure and
transport steps conditional on those identities. It does not
claim to have proved those analytic identities from formal
definitions alone.

\paragraph{TL19: finite generators for all bounded-degree normalized
jets.}
Assume the actual smoothness, unitary covariance, normalization
character law, and \texttt{JetAlg} of the theta family in question;
also assume its $c_D$ are algebraic, as in TL17. Fix positive
modulus $M$, a finite prime set, two ranks, and a total degree
$d$. TL1 gives finitely many self-dual lattice representatives
in the two ranks. There are finitely many divisors $c\mid M$.
For fixed representatives and $c$, the permitted rational shift
is bounded in a fixed fractional Hom lattice, while the theta
function is unchanged by adding $c$ times its local integral
Hom lattice. The resulting quotient is finite: localize at the
finitely many relevant primes, where a bounded lattice modulo
a deeper full lattice has finite residue cardinality, and use
the injective map to those local quotients as in TL2.
There are also finitely many multi-indices of degree at most
$d$.

Take all normalized derivatives for these finite choices.
They are algebraic by \texttt{JetAlg}. Adjoin them, the finitely
many $c_D$, and generators for all embedded copies of $E$.
For arbitrary self-dual lattices choose rational unitary maps
to the fixed representatives. The actual covariance identity
expresses the theta function as the representative function
composed with its specified linear change of variables, whose
coefficients are embedded elements of $E$. Repeated chain rule
expresses each derivative of degree $d$ as a finite sum of
representative derivatives of that degree times products of
these coefficients. The normalization-character change is
the prescribed product of embedded determinants, also in the
same field. Finally the rescaling by $c_D$ uses algebraic
factors already adjoined. This proves \texttt{FinAlg(jetSet)}:
one finite algebraic extension contains the whole infinite
set of these jets.

For canonical column data the exact modulus is
\[
 M(u,S_G)=Q^{m+1}\bigl(T_u Q^{m+1}\bigr),
 \qquad Q=\prod_{P\in S_G}q_P.                    \tag{TL.21}
\]
All prime sets used by the steps are subsets of the finite
$\texttt{badP}(B,\Xi)$ for the finite character family $\Xi$.
There are finitely many words, steps $i<m$, and such pairs of
prime sets. The union of their \texttt{FinAlg} sets is again
\texttt{FinAlg}, by taking the finite union of generators.
The source's \texttt{allSet} includes only indices satisfying
\texttt{FinAlg}; its finiteness proof alone does not show that
an actual step belongs to it. The analytic inputs and the
preceding argument supply exactly that missing membership.
The actual multidegrees have total degree \texttt{totDeg},
so containment of the chosen generators in the coefficient
field proves the literal \texttt{JetsIn} for those canonical
data. This is a finite-generation reduction, not a proof of
analytic algebraicity by selecting an empty set.

\paragraph{TL20: what the local index and canonicalization really
require.}
At an ordinary twist prime the initial index construction selects
finite representatives recursively. Starting with $F_0=\{1\}$,
retain at the next step exactly the admitted double cosets which
have a nonzero first-unfolding coefficient from a member of $F_i$.
Given finite admitted level-coset representatives, normal-form
reduction, and invariance of nonvanishing under the indicated
unipotent/double-coset moves, choose one normalized representative
in each retained coset. This produces a finite $F_{i+1}$, with
the exact predecessor, uniqueness, and exhaustivity properties.
At $m-1$ set $F_{m-1}=\{1\}$ and above it set $F_j=\varnothing$,
as stipulated by the source. Finite products over the local
primes give the global index sets. A level character descends
to its projected group when it is trivial on the kernel:
define its value using a lift, and compare two lifts by their
quotient in that kernel. This explains the finite selection
and character-descent parts of \texttt{exists\_chIdx}; the
needed local coset, normal-form, kernel-character, and
first-unfolding theorems remain genuine inputs here.

For an already supplied $I.\texttt{Ok}$ of depth at least $2m$,
much of the canonical arithmetic can be extracted without an
existence argument. Its square-root clause gives $\sqrt{q_P}$.
Its reference local-character formula gives the required
coefficient-field values of $\mu_j(q_P)^{-1}$. If the old depth
is $D=2m+a$, evaluate its Weil identity at
$\xi_0^{-1}/q_P^{2m}$, using the integer $q_P^a$ in that
identity. The local-unit clause makes this an allowed argument
and gives exactly $\exp(2\pi\mathrm i/q_P^{2m})$ in the
coefficient field. Thus the primitive roots required for the
canonical depth are already present. The base congruences give
$q_P\equiv1\pmod4$ at auxiliary primes and the tuple separates
twist primes from them.

Replacing the constants and indices by \texttt{twCst} and
\texttt{twIdx} now produces \texttt{IsTw} definitionally,
and leaves the host, character family, tuple readings, selection,
and ambient theta data unchanged. To assert that the replacement
still satisfies \texttt{Ok}, however, the source invokes the
actual twisted first-unfolding nonvanishing and its double-coset
invariance, together with the corresponding character
nonvanishing theorem. These do not follow from the root and
square-root extraction just proved. They remain the exact
\texttt{twIdx\_ok\_closed} boundary. Canonical replacement
therefore does not make the local-index hypotheses disappear.

\paragraph{TL21: the quantified canonical evaluation assembly and
the remaining construction boundary.}
The phrase \texttt{EvLg} has two nested meanings. Its inner
\texttt{large\_2} assertion chooses finite algebraic generator
sets $S_1$ and $S_2$, and finite prime/avoidance conditions,
before quantifying over a base $B$ lying over the fixed context,
a frame $f$, and a twist datum $T$. The antecedent requires
$S_1$ in the \emph{fixed} coefficient field and $S_2$ in the full
coefficient field. It does not assert the result for an arbitrary
unchanged coefficient field. In this constructor the additional
prime and avoidance sets can be empty. The twist family
\[
 \Xi=\{1,\delta,\delta'\}\cup
       \{(\delta'\text{ if }\operatorname{reschar}(w)=p_\delta
                      \text{ else }\delta)\psi_w:w\}
 \tag{TL.22}
\]
is finite: outside the finite set of supported rational primes
$\psi_w=1$, and there are only finitely many coefficient places
over each supported prime.

Choose $S_1$ to contain a primitive generator for TL14's common
frame-value field. Choose $S_2$ to contain the period-field
generators of TL17, the roots and square roots of TL15, and the
jet generators of TL19. Their finite union is algebraic under
the primitive inputs explicitly retained in TL17--TL20.
With a nonempty \texttt{PreX}, TL16 gives the required auxiliary
orientation and type-prime quotient identity. Under the stated
frame-inclusion antecedent this gives both fields of
\texttt{Pins}.

The finite character set $\Xi$ itself is the index type. TL15
gives its unique tuples, the three distinguished indices, the
trivial reference tuple, and the actual residue character data.
The local-index suppliers in TL20, applied with initial
$M=A=B=m$ and twist depth $2m$, then give the actual initial
\texttt{TC.Ok}; the supplied ambient period/trace-isometry
constructor supplies its specified \texttt{AmbIn}. The selection,
frame, excluded set, and range of $\xi$ are exactly the prescribed
ones, because they were used as these record fields. TL20's
canonical index supplier preserves \texttt{Ok} and gives
\texttt{IsTw}; TL19 gives \texttt{JetsIn}. TL14 supplies
\texttt{ValLL0} for every later assembly with the same base and
frame, and TL17 supplies \texttt{FieldOK}. Finally TL13 now
supplies \texttt{Laws} without a further premise.

The outer \texttt{Ev} quantifies a positive $\varepsilon_0$, then
$\varepsilon,C_8$, a transported lower-bound datum, constants,
and sufficiently large contexts. The inner construction just
described works for every such context, so the outer thresholds
can be chosen trivially (for example $\varepsilon_0=1$ and
zero size thresholds) while retaining the literal inner
coefficient-field requirements. This explains the source's
\texttt{ev\_of\_forall} step; it neither constructs an arithmetic
host nor removes the existing context hypotheses.

Accordingly, the remaining sources for the full canonical
existence conclusion are the actual local-index/nonvanishing
theorems; the period/trace-isometry ambient datum and CM
$j$-algebraicity; and the primitive normalized-theta identities
specified in TL18 together with their unitary covariance law.
The finite fields, coefficient roots, tuple reading, orientation,
ideal principalization, finite derivative reduction, and the
entire implication \texttt{Ok}$\Rightarrow$\texttt{Laws} have
been reconstructed above. The record \texttt{Later} still
consists of five separately supplied stages. A universal
statement about all such later assemblies is not an existence
proof of those stages, of a primitive host, or of RH/GRH.

